\chapter{Kesimpulan dan Saran}

\section{Kesimpulan}

Penelitian ini merancang dan mengimplementasikan \textbf{IP core transmitter I2S} dengan antarmuka \textbf{AXI4-Lite} pada SoC \textbf{MicroBlaze V}, board \textbf{Basys3}, dan keluaran menuju \textbf{PCM5102A}. Berdasarkan hasil pada Bab IV, dapat disimpulkan:

\begin{enumerate}
	\item IP memungkinkan kontrol audio melalui register memory-mapped sesuai register map yang ditargetkan, termasuk jalur \textbf{FIFO TX} dan \textbf{interupsi watermark} untuk mengurangi underrun.
	\item Dua mode frekuensi sampling (\textbf{sekitar} 48 kHz dan 96 kHz) diperoleh dari pembagi \texttt{audio\_clk}; nilai terukur mengikuti frekuensi \textit{actual} Clocking Wizard dan harus dilaporkan secara eksplisit.
	\item Verifikasi dengan \textbf{ILA} dan bukti audio menunjukkan serializer \textbf{Philips I2S} 24-bit dalam slot 32-bit bekerja pada perangkat keras yang diuji.
\end{enumerate}

\section{Saran}

\begin{enumerate}
	\item Menambahkan \textbf{DMA} untuk mengurangi beban CPU saat streaming panjang.
	\item Mendukung \textbf{RX} atau format \textbf{TDM}/multi-channel untuk aplikasi mikrofon array.
	\item Menyusun \textbf{device tree} dan driver (mis.\ ASoC Linux) bila IP akan dipakai pada platform ber-OS.
	\item Menyempurnakan \textbf{CDC} antara domain AXI dan audio jika FIFO dan serializer berada pada domain clock berbeda dengan metode resynchronisasi formal.
	\item Menguji pada \textbf{osilator referensi} eksternal atau nilai MMCM yang lebih mendekati 24,576 MHz untuk mendekati \(F_s\) nominal persis.
\end{enumerate}

\noindent\textit{Penulis disarankan menyesuaikan poin kesimpulan dengan data terukur final dan arahan pembimbing.}
